\documentclass[a4paper,12pt,twoside]{memoir}

% Castellano
\usepackage[spanish,es-tabla]{babel}
\selectlanguage{spanish}
\usepackage[utf8]{inputenc}
\usepackage[T1]{fontenc}
\usepackage{lmodern} % scalable font
\usepackage{microtype}
\usepackage{placeins}

\RequirePackage{booktabs}
\RequirePackage[table]{xcolor}
\RequirePackage{xtab}
\RequirePackage{multirow}

% Links
\PassOptionsToPackage{hyphens}{url}\usepackage[colorlinks]{hyperref}
\hypersetup{
	allcolors = {red}
}

% Ecuaciones
\usepackage{amsmath}

% Rutas de fichero / paquete
\newcommand{\ruta}[1]{{\sffamily #1}}

% Párrafos
\nonzeroparskip

% Huérfanas y viudas
\widowpenalty100000
\clubpenalty100000

% Evitar solapes en el header
\nouppercaseheads

% Cuadros de notas (comando propuesto por la alumna)
% Los parámetros son:
% 1 --> Texto de la nota
\newcommand{\nota}[1]{
    \begin{center}
        \fcolorbox{blue}{white}{\textcolor{black}{
            \begin{minipage}{0.915\linewidth}
        	\textbf{\color{blue}NOTA}
                \linebreak\linebreak
                #1
                \linebreak
            \end{minipage}
        }}
    \end{center}
}

% Imagenes
\usepackage{graphicx}
\newcommand{\imagen}[3]{
	\begin{figure}[!h]
		\centering
		\includegraphics[width=#3\textwidth]{#1}
		\caption{#2}\label{fig:#1}
	\end{figure}
	\FloatBarrier
}
\newcommand{\imagenflotante}[2]{
	\begin{figure}%[!h]
		\centering
		\includegraphics[width=0.9\textwidth]{#1}
		\caption{#2}\label{fig:#1}
	\end{figure}
}



% El comando \figura nos permite insertar figuras comodamente, y utilizando
% siempre el mismo formato. Los parametros son:
% 1 -> Porcentaje del ancho de página que ocupará la figura (de 0 a 1)
% 2 --> Fichero de la imagen
% 3 --> Texto a pie de imagen
% 4 --> Etiqueta (label) para referencias
% 5 --> Opciones que queramos pasarle al \includegraphics
% 6 --> Opciones de posicionamiento a pasarle a \begin{figure}
\newcommand{\figuraConPosicion}[6]{%
  \setlength{\anchoFloat}{#1\textwidth}%
  \addtolength{\anchoFloat}{-4\fboxsep}%
  \setlength{\anchoFigura}{\anchoFloat}%
  \begin{figure}[#6]
    \begin{center}%
      \Ovalbox{%
        \begin{minipage}{\anchoFloat}%
          \begin{center}%
            \includegraphics[width=\anchoFigura,#5]{#2}%
            \caption{#3}%
            \label{#4}%
          \end{center}%
        \end{minipage}
      }%
    \end{center}%
  \end{figure}%
}

%
% Comando para incluir imágenes en formato apaisado (sin marco).
\newcommand{\figuraApaisadaSinMarco}[5]{%
  \begin{figure}%
    \begin{center}%
    \includegraphics[angle=90,height=#1\textheight,#5]{#2}%
    \caption{#3}%
    \label{#4}%
    \end{center}%
  \end{figure}%
}
% Para las tablas
\newcommand{\otoprule}{\midrule [\heavyrulewidth]}
%
% Nuevo comando para tablas pequeñas (menos de una página).
\newcommand{\tablaSmall}[5]{%
 \begin{table}
  \begin{center}
   \rowcolors {2}{gray!35}{}
   \begin{tabular}{#2}
    \toprule
    #4
    \otoprule
    #5
    \bottomrule
   \end{tabular}
   \caption{#1}
   \label{tabla:#3}
  \end{center}
 \end{table}
}

%
%Para el float H de tablaSmallSinColores
\usepackage{float}

%
% Nuevo comando para tablas pequeñas (menos de una página).
\newcommand{\tablaSmallSinColores}[5]{%
 \begin{table}[H]
  \begin{center}
   \begin{tabular}{#2}
    \toprule
    #4
    \otoprule
    #5
    \bottomrule
   \end{tabular}
   \caption{#1}
   \label{tabla:#3}
  \end{center}
 \end{table}
}

\newcommand{\tablaApaisadaSmall}[5]{%
\begin{landscape}
  \begin{table}
   \begin{center}
    \rowcolors {2}{gray!35}{}
    \begin{tabular}{#2}
     \toprule
     #4
     \otoprule
     #5
     \bottomrule
    \end{tabular}
    \caption{#1}
    \label{tabla:#3}
   \end{center}
  \end{table}
\end{landscape}
}

%
% Nuevo comando para tablas grandes con cabecera y filas alternas coloreadas en gris.
\newcommand{\tabla}[6]{%
  \begin{center}
    \tablefirsthead{
      \toprule
      #5
      \otoprule
    }
    \tablehead{
      \multicolumn{#3}{l}{\small}\\
      \toprule
      #5
      \otoprule
    }
    \tabletail{
      \hline
      \multicolumn{#3}{r}{\small}\\
    }
    \tablelasttail{
      \hline
    }
    \bottomcaption{#1}
    \rowcolors {2}{gray!35}{}
    \begin{xtabular}{#2}
      #6
      \bottomrule
    \end{xtabular}
    \label{tabla:#4}
  \end{center}
}
\newcommand{\tablatres}[6]{%
  \begin{center}
    \tablefirsthead{
      \toprule
      #5
      \otoprule
    }
    \tablehead{
      \multicolumn{#3}{l}{\small}\\
      \toprule
      #5
      \otoprule
    }
    \tabletail{
      \hline
      \multicolumn{#3}{r}{\small}\\
    }
    \tablelasttail{
      \hline
    }
    %\bottomcaption{#1}
    \rowcolors {2}{gray!35}{}
    \begin{xtabular}{#2}
      #6
      \bottomrule
    \end{xtabular}
  \end{center}
}

%
% Nuevo comando para tablas grandes con cabecera.
\newcommand{\tablaSinColores}[6]{%
  \begin{center}
    \tablefirsthead{
      \toprule
      #5
      \otoprule
    }
    \tablehead{
      \multicolumn{#3}{l}{\small\sl continúa desde la página anterior}\\
      \toprule
      #5
      \otoprule
    }
    \tabletail{
      \hline
      \multicolumn{#3}{r}{\small\sl continúa en la página siguiente}\\
    }
    \tablelasttail{
      \hline
    }
    \bottomcaption{#1}
    \begin{xtabular}{#2}
      #6
      \bottomrule
    \end{xtabular}
    \label{tabla:#4}
  \end{center}
}

%
% Nuevo comando para tablas grandes sin cabecera.
\newcommand{\tablaSinCabecera}[5]{%
  \begin{center}
    \tablefirsthead{
      \toprule
    }
    \tablehead{
      \multicolumn{#3}{l}{\small\sl continúa desde la página anterior}\\
      \hline
    }
    \tabletail{
      \hline
      \multicolumn{#3}{r}{\small\sl continúa en la página siguiente}\\
    }
    \tablelasttail{
      \hline
    }
    \bottomcaption{#1}
  \begin{xtabular}{#2}
    #5
   \bottomrule
  \end{xtabular}
  \label{tabla:#4}
  \end{center}
}



\definecolor{cgoLight}{HTML}{EEEEEE}
\definecolor{cgoExtralight}{HTML}{FFFFFF}

%
% Nuevo comando para tablas grandes sin cabecera.
\newcommand{\tablaSinCabeceraConBandas}[5]{%
  \begin{center}
    \tablefirsthead{
      \toprule
    }
    \tablehead{
      \multicolumn{#3}{l}{\small\sl continúa desde la página anterior}\\
      \hline
    }
    \tabletail{
      \hline
      \multicolumn{#3}{r}{\small\sl continúa en la página siguiente}\\
    }
    \tablelasttail{
      \hline
    }
    \bottomcaption{#1}
    \rowcolors[]{1}{cgoExtralight}{cgoLight}

  \begin{xtabular}{#2}
    #5
   \bottomrule
  \end{xtabular}
  \label{tabla:#4}
  \end{center}
}




\graphicspath{ {./img/} }

% Capítulos
\chapterstyle{bianchi}
\newcommand{\capitulo}[2]{
	\setcounter{chapter}{#1}
	\setcounter{section}{0}
	\setcounter{figure}{0}
	\setcounter{table}{0}
	\chapter*{#2}
	\addcontentsline{toc}{chapter}{#2}
	\markboth{#2}{#2}
}

% Apéndices
\renewcommand{\appendixname}{Apéndice}
\renewcommand*\cftappendixname{\appendixname}

\newcommand{\apendice}[1]{
	%\renewcommand{\thechapter}{A}
	\chapter{#1}
}

\renewcommand*\cftappendixname{\appendixname\ }

% Formato de portada
\makeatletter
\usepackage{xcolor}
\newcommand{\tutor}[1]{\def\@tutor{#1}}
\newcommand{\course}[1]{\def\@course{#1}}
\definecolor{cpardoBox}{HTML}{E6E6FF}
\def\maketitle{
  \null
  \thispagestyle{empty}
  % Cabecera ----------------
\noindent\includegraphics[width=\textwidth]{cabecera}\vspace{1cm}%
  \vfill
  % Título proyecto y escudo informática ----------------
  \colorbox{cpardoBox}{%
    \begin{minipage}{.8\textwidth}
      \vspace{.5cm}\Large
      \begin{center}
      \textbf{TFG del Grado en Ingeniería Informática}\vspace{.6cm}\\
      \textbf{\LARGE\@title{}}
      \end{center}
      \vspace{.2cm}
    \end{minipage}

  }%
  \hfill\begin{minipage}{.20\textwidth}
    \includegraphics[width=\textwidth]{escudoInfor}
  \end{minipage}
  \vfill
  % Datos de alumno, curso y tutores ------------------
  \begin{center}%
  {%
    \noindent\LARGE
    Presentado por \@author{}\\ 
    en Universidad de Burgos --- \@date{}\\
    Tutores: \@tutor{}\\
  }%
  \end{center}%
  \null
  \cleardoublepage
  }
\makeatother


% Datos de portada
\title{Impact Factor Oracle\\Documentación Técnica}
\author{Gadea Lucas Pérez}
\tutor{Virginia Ahedo y Álvar Arnaiz}
\date{\today}

\begin{document}

\maketitle



\cleardoublepage



%%%%%%%%%%%%%%%%%%%%%%%%%%%%%%%%%%%%%%%%%%%%%%%%%%%%%%%%%%%%%%%%%%%%%%%%%%%%%%%%%%%%%%%%



\frontmatter


\clearpage

% Indices
\tableofcontents

\clearpage

\listoffigures

\clearpage

\listoftables

\clearpage

\mainmatter

\appendix


\apendice{Plan de Proyecto Software}

\section{Introducción}

La planificación temporal es esencial para el éxito de cualquier proyecto, especialmente en el desarrollo de \textit{software}. En esta sección se detallará cómo se ha llevado a cabo el cronograma siguiendo una metodología ágil tipo Scrum, mediante el uso de \textit{sprints}. Se describirán los pasos necesarios para planificar y gestionar de manera eficiente el tiempo y los recursos disponibles, asegurando así el cumplimiento de los objetivos del proyecto en el plazo establecido.

\section{Metodología}
La metodología Scrum es un marco de trabajo ágil que se utiliza para la gestión de proyectos. Se basa en la colaboración entre el equipo de desarrollo y el cliente, así como en la entrega continua de entregables funcionales en ciclos cortos de tiempo conocidos como \textit{sprints}. Scrum se centra en la flexibilidad y la adaptabilidad, permitiendo más fácilmente la adaptación y los cambios de manera rápida y efectiva ante situaciones inesperadas. 

En el contexto del presente trabajo, el uso de la metodología Scrum será beneficioso puesto que nos permitirá tener un enfoque más dinámico y flexible, lo que nos permitirá adaptarnos a las necesidades del proyecto a medida que avanzamos. Además, Scrum nos proporcionará una mayor transparencia y comunicación, lo que nos permitirá tomar decisiones informadas y asegurar que el proyecto se entregue a tiempo y con éxito.

Por otro lado, se ha complementado la metodología Scrum con el modelo Canvas\footnote{El modelo Canvas se refiere comúnmente a un lienzo o plantilla visual utilizada para estructurar y desarrollar ideas de negocio.}. Esta combinación nos permitirá visualizar y organizar de manera efectiva las tareas, los hitos y los recursos disponibles, facilitando así la gestión y el seguimiento del proyecto.

\section{Herramienta}
ZenHub es una herramienta de gestión de proyectos que se utiliza en muchos entornos empresariales para mejorar la eficiencia en el seguimiento y administración de tareas. Con ZenHub, es posible planificar y visualizar fácilmente las actividades de un proyecto, monitorear el progreso de cada tarea y hacer un seguimiento de los plazos de entrega. También ofrece funciones útiles como integración con GitHub, herramientas de informes y análisis de datos, y seguimiento de errores y problemas.

\section{Planificación temporal}

El presente proyecto comienza en septiembre y se extiende hasta junio. Durante este período de tiempo, es importante asegurarnos de que todas las tareas y hitos estén claramente definidos.

Para lograr esto, se han realizado reuniones periódicas para discutir el progreso del proyecto y asegurarnos de que estamos en el camino correcto. Además, hemos implementado \textit{sprints} regulares para asegurarnos de que estamos avanzando de manera constante y cumpliendo con nuestras metas a tiempo.

Con esta planificación temporal sólida, estamos seguros de que podremos completar el proyecto a tiempo y cumplir con los objetivos establecidos. Sin embargo, es importante ser flexibles y estar preparados para hacer ajustes según sea necesario a medida que avanzamos en el proyecto, puesto que se realiza al mismo tiempo que trascurre el curso académico.

En las siguientes tablas se recogen los distintos \textit{sprints} junto con su duración, objetivos y tareas. Además, se adjuntará el gráfico \textit{burndown}\footnote{Gráfico burndown: Se trata de una herramienta comúnmente utilizada en la gestión de proyectos ágiles para realizar un seguimiento de un \textit{sprint}. Este gráfico muestra la cantidad de trabajo que queda por hacer en relación con el tiempo disponible para hacerlo. A medida que se completa el trabajo, la línea del gráfico debería disminuir hasta alcanzar cero en el momento en que finaliza el \textit{sprint}.} de cada uno de los \textit{sprints}.


\tabla{Sprint 1}{llll}{4}{sprint1}{\textbf{Sprt.} & \textbf{Duración} & \textbf{Objetivos} & \textbf{Tareas}\\}{
    1 & \parbox{55}{19/09/2022 03/10/2022} & \parbox{80}{Comenzar el desarrollo del proyecto} & \parbox{150}{\begin{itemize}\item Elegir un modelo de referencias bibliográficas. \item Definición de conceptos.\end{itemize}}\\
}

\imagen{sprint1}{Sprint 1: Burndown chart}{0.9}

\tabla{Sprint 2}{llll}{4}{sprint2}{\textbf{Sprt.} & \textbf{Duración} & \textbf{Objetivos} & \textbf{Tareas}\\}{
    2 & \parbox{55}{03/10/2022 17/10/2022} & \parbox{80}{Tareas de investigación y documentación} & \parbox{150}{\begin{itemize}\item Búsqueda de antecedentes. \item Familiarizarse con la memoria en \LaTeX.\end{itemize}}\\
}

\imagen{sprint2}{Sprint 2: Burndown chart}{0.9}

\tabla{Sprint 3}{llll}{4}{sprint3}{\textbf{Sprt.} & \textbf{Duración} & \textbf{Objetivos} & \textbf{Tareas}\\}{
    3 & \parbox{55}{17/10/2022 14/11/2022} & \parbox{80}{Investigación y creación de un prototipo inicial} & \parbox{150}{\begin{itemize}\item Búsqueda de información sobre MIAR. \item Documentación sobre el JCR y el índice de impacto. \item Prototipo que permita la búsqueda de artículos en GS. \item Extracción de datos: ¿qué datos se pueden extraer? \item Reflexión sobre el diseño de la BBDD. \item Prueba del prototipo.\end{itemize}}\\
}

\imagen{sprint3}{Sprint 3: Burndown chart}{0.9}

\tabla{Sprint 4}{llll}{4}{sprint4}{\textbf{Sprt.} & \textbf{Duración} & \textbf{Objetivos} & \textbf{Tareas}\\}{
    4 & \parbox{55}{14/11/2022 28/11/2022} & \parbox{80}{Base de datos y prototipo} & \parbox{150}{\begin{itemize}\item Creación de la base de datos en MariaDB. \item Prueba de la BBDD. \item Mejoras del prototipo.\end{itemize}}\\
}

\imagen{sprint4}{Sprint 4: Burndown chart}{0.85}

\nota{Como se puede observar en el gráfico \textit{burndown} anterior, las tareas no fueron cerradas a tiempo. Sin embargo, es preciso aclarar que se trata de un error en la configuración de las tareas. Si bien todas las tareas fueron completadas a tiempo, fue preciso borrarlas en ZenHub y volverlas a crear más adelante cuando el \textit{sprint} ya había terminado debido a que estaban inicialmente mal configuradas.}

\tabla{Sprint 5}{llll}{4}{sprint5}{\textbf{Sprt.} & \textbf{Duración} & \textbf{Objetivos} & \textbf{Tareas}\\}{
    5 & \parbox{55}{28/11/2022 12/12/2022} & \parbox{80}{Base de datos y prototipo} & \parbox{150}{\begin{itemize}\item Se muda la BBDD a Postgres. \item Pruebas e investigación para extraer el DOI.\end{itemize}}\\
}

\imagen{sprint5}{Sprint 5: Burndown chart}{0.9}

\tabla{Sprint 6}{llll}{4}{sprint6}{\textbf{Sprt.} & \textbf{Duración} & \textbf{Objetivos} & \textbf{Tareas}\\}{
    6 & \parbox{55}{09/01/2023 23/01/2023} & \parbox{80}{Documentación y obtención del DOI} & \parbox{150}{\begin{itemize}\item Obtención del DOI a través de Crossref. \item Avance de la memoria y los anexos. \item Mejora de la BBDD.\end{itemize}}\\
}

\imagen{sprint6}{Sprint 6: Burndown chart}{0.9}

\tabla{Sprint 7}{llll}{4}{sprint7}{\textbf{Sprt.} & \textbf{Duración} & \textbf{Objetivos} & \textbf{Tareas}\\}{
    7 & \parbox{55}{06/02/2023 20/02/2023} & \parbox{80}{Preparación para la conexión remota a los servidores de la universidad.} & \parbox{150}{\begin{itemize}\item Conectarse a la VPN de la universidad. \item Crear un servicio virtualizado en Miniconda. \item Mejorar la BBDD.\end{itemize}}\\
}

\imagen{sprint7}{Sprint 7: Burndown chart}{0.9}

\tabla{Sprint 8}{llll}{4}{sprint8}{\textbf{Sprt.} & \textbf{Duración} & \textbf{Objetivos} & \textbf{Tareas}\\}{
    8 & \parbox{55}{20/02/2023 06/03/2023} & \parbox{80}{Conexión a los servidores y mejora de los prototipos. Avances en la API y documentación.} & \parbox{150}{\begin{itemize}\item Mejorar el prototipo. \item Prototipo con varios hilos de ejecución (usar servidores de la universidad). \item Nuevo prototipo con Selenium. \item Sistema inteligente de requests. \item Revisión de PoP. \item API: Autentificación y back-end.\end{itemize}}\\
}

\imagen{sprint8}{Sprint 8: Burndown chart}{0.9}

\tabla{Sprint 9}{llll}{4}{sprint9}{\textbf{Sprt.} & \textbf{Duración} & \textbf{Objetivos} & \textbf{Tareas}\\}{
    9 & \parbox{55}{06/03/2023 20/03/2023} & \parbox{80}{Finalización del prototipo de extracción de Crossref.} & \parbox{150}{\begin{itemize}\item Plan de prueba para Crossref. \item Prototipo de Crossref. \item Comprobar límite de llamadas a Crossref. \item Crear entorno virtualizado. \item Fichero requirements. \item Licencia Crossref para Cited-by. \item Cálculo del IF.\end{itemize}}\\
}

\imagen{sprint9}{Sprint 9: Burndown chart}{0.9}

\tabla{Sprint 10}{llll}{4}{sprint10}{\textbf{Sprt.} & \textbf{Duración} & \textbf{Objetivos} & \textbf{Tareas}\\}{
    10 & \parbox{55}{20/03/2023 03/04/2023} & \parbox{80}{Procesamiento de los datos.} & \parbox{150}{\begin{itemize}\item Normalizar listados JCR. \item Comparar JCR obtenidos con los valores reales.\end{itemize}}\\
}


Finalmente, se incluirá un gráfico acumulativo\footnote{Gráfico acumulativo: En este gráfico, el eje vertical representa la cantidad total de trabajo completado, mientras que el eje horizontal representa el tiempo transcurrido durante el proyecto. A medida que se completa el trabajo, la línea del gráfico se eleva, lo que indica el progreso general del proyecto} que muestra la cantidad total de trabajo realizado de forma más generalizada.

\imagen{cumulative}{Gráfico acumulativo de todo el proyecto}{0.9}




\section{Estudio de viabilidad}

\subsection{Viabilidad económica}

\subsection{Viabilidad legal}



\apendice{Especificación de Requisitos}

\section{Introducción}

Para llevar a cabo el proyecto adecuadamente, se realizará un análisis exhaustivo de los requisitos funcionales y no funcionales de la aplicación. En la sección de Objetivos generales (\ref{sec:Objetivos generales}) se establecerán los resultados esperados y los criterios de éxito del proyecto. Seguidamente, se presentará el catálogo de requisitos (sección \ref{sec:Catálogo de requisitos}), detallando las funcionalidades y características que la aplicación debe cumplir. Finalmente, se procederá a la especificación de requisitos (sección \ref{sec:Especificación de requisitos}), describiendo cada uno de los identificados en el catálogo y estableciendo su prioridad, complejidad y dependencias.

A través de este informe, se pretende sentar las bases y establecer las directrices necesarias para el desarrollo exitoso de la aplicación web que proporcionará a los usuarios la información clave sobre el impacto de las revistas científicas, con el objetivo de facilitar la toma de decisiones informadas en el ámbito de la publicación académica.

\section{Objetivos generales}
\label{sec:Objetivos generales}

Los objetivos generales de este proyecto son los siguientes:

\begin{enumerate}
    \item Realizar el proceso de extracción de datos relevantes para el cálculo del Índice de Impacto de las revistas científicas. Esta tarea constituye la parte más desafiante del proyecto, ya que implica llevar a cabo una investigación exhaustiva y realizar pruebas utilizando diversas técnicas de \textit{web scraping} y simulación con Selenium. Se requerirá desarrollar métodos eficientes para recopilar los datos necesarios que no estén fácilmente accesibles y estructurados.
    
    \item Implementar modelos de predicción utilizando técnicas de \textit{machine learning)}, en particular utilizando la biblioteca Scikit-learn. Estos modelos permitirán predecir el Índice de Impacto de las revistas científicas con base en las características y métricas disponibles. Se explorarán diferentes algoritmos y técnicas para obtener predicciones precisas y confiables.
    
    \item Desarrollar una aplicación web utilizando el \textit{framework} Flask. Esta aplicación será la interfaz principal para que los usuarios puedan acceder y visualizar los resultados del proyecto. La aplicación web proporcionará una experiencia de usuario intuitiva y amigable, permitiendo la búsqueda y filtrado de revistas, así como la visualización de los índices de impacto calculados y las predicciones generadas por los modelos.
\end{enumerate}

El objetivo central del proyecto es proporcionar a los usuarios una herramienta integral que les permita evaluar y comparar el impacto de las revistas científicas de manera eficiente y precisa. Para lograr esto, se abordarán tres aspectos clave: la extracción de datos, los modelos de predicción y el desarrollo de la aplicación web. Cada uno de estos objetivos se complementa para lograr un sistema funcional y útil para la comunidad académica y científica.

Es importante reiterar que el proceso de extracción de datos es considerado como la parte más desafiante del proyecto debido a la necesidad de investigar y probar diferentes técnicas de \textit{web scraping}. El éxito en esta etapa es fundamental para garantizar la disponibilidad de los datos necesarios para el cálculo del Índice de Impacto y para el entrenamiento de los modelos de predicción.


\section{Catálogo de Historias de Usuario}
\label{sec:Catálogo de requisitos}

En este caso, dado que se ha elegido seguir una metodología ágil, se considera un catálogo de historias de usuario (usadas en ambos marcos, Scrum y Kanban). Para poder redactar una historia de usuario, hay tres elementos clave~\cite{Asana_2022}:
\begin{itemize}
    \item Perfil: El rol del usuario final.
    \item Necesidad: El objetivo que tiene la función de \textit{software} para el usuario final.
    \item Propósito: El objetivo de la experiencia del usuario final con la función de \textit{software}.
\end{itemize}

Habiendo identificado estos elementos, la historia de usuario se redactará siguiendo este formato: \textbf{<<Como [perfil], quiero [necesidad], para lograr [propósito]>>}.

Se ha elaborado el catálogo de historias de usuario con el objetivo de cubrir todos los objetivos mencionados en el apartado anterior sin entrar en detalles (los requisitos más específicos con respecto a la aplicación web están definidos en la sección \ref{sec:Requisitos de aplicación}). Para visualizar mejor el catálogo, se han clasificado las historias de usuario en tres categorías: extracción de datos (Tabla \ref{tab:extracción_de_datos}), modelos de predicción  (Tabla \ref{tab:Modelos de predicción}) y aplicación web (Tabla \ref{tab:Aplicación web}).

\begin{table}[h]
\centering
\begin{tabular}{|p{3cm}|p{8cm}|}
\hline
\textbf{Identificador} & \textbf{Descripción} \\
\hline
HU1 & Como desarrollador, quiero obtener información relevante para el cálculo del Índice de Impacto (JCR) de las revistas científicas a partir de fuentes externas. \\
\hline
HU2 & Como desarrollador, quiero aplicar técnicas de \textit{web scraping} y simulación con Selenium para extraer los datos de manera eficiente y precisa. \\
\hline
HU3 & Como desarrollador, quiero procesar y almacenar los datos extraídos en una base de datos para su posterior uso en el cálculo del Índice de Impacto y en los modelos de predicción. \\
\hline
\end{tabular}
\caption{HU Extracción de datos}
\label{tab:extracción_de_datos}
\end{table}
\newpage

\begin{table}[h]
\centering
\begin{tabular}{|p{3cm}|p{8cm}|}
\hline
\textbf{Identificador} & \textbf{Descripción} \\
\hline
HU4 & Como desarrollador, quiero implementar modelos de predicción utilizando técnicas de \textit{machine learning} (scikit-learn) para predecir el Índice de Impacto de las revistas científicas. \\
\hline
HU5 & Como desarrollador, quiero evaluar y comparar diferentes algoritmos de \textit{machine learning} para seleccionar los modelos más adecuados. \\
\hline
HU6 & Como desarrollador, quiero entrenar los modelos utilizando los datos almacenados y ajustar sus parámetros para obtener predicciones precisas. \\
\hline
\end{tabular}
\caption{HU Modelos de predicción}
\label{tab:Modelos de predicción}
\end{table}

\begin{table}[h]
\centering
\begin{tabular}{|p{3cm}|p{8cm}|}
\hline
\textbf{Identificador} & \textbf{Descripción} \\
\hline
HU7 & Como usuario, quiero acceder a una aplicación web donde pueda buscar y filtrar revistas científicas para evaluar su Índice de Impacto. \\
\hline
HU8 & Como usuario, quiero visualizar de forma clara y comprensible los índices de impacto calculados y las predicciones generadas por los modelos de predicción. \\
\hline
HU9 & Como usuario, quiero interactuar con la aplicación web, seleccionar revistas específicas y obtener información detallada sobre cada una de ellas. \\
\hline
\end{tabular}
\caption{HU Aplicación web}
\label{tab:Aplicación web}
\end{table}

Estas historias de usuario representan las necesidades y objetivos de los diferentes perfiles de usuarios en relación con la extracción de datos, los modelos de predicción y la aplicación web. Siguiendo una metodología ágil, estas historias de usuario servirán como guía para el desarrollo incremental y la entrega de funcionalidades de valor en cada iteración del proyecto.

\newpage
\section{Requisitos de la aplicación}
\label{sec:Requisitos de aplicación}

En esta sección se enumerarán, en mayor detalle, los requisitos que se esperan de la aplicación web en concreto.

\subsection{Requisitos funcionales}

\begin{itemize}
  \item \textbf{RF1}: Autenticación y Registro
  \begin{itemize}
    \item \textbf{RF1.1}: La aplicación debe permitir a los usuarios autenticarse mediante un nombre de usuario y contraseña válidos.
    \item \textbf{RF1.2}: La aplicación debe permitir a los usuarios registrarse proporcionando un nombre de usuario, una contraseña, una dirección de correo electrónico válida y confirmación de la contraseña.
    \item \textbf{RF1.3}: Las contraseñas deben tener un formato válido (i.e.: una longitud mínima de 8 caracteres, caracteres especiales requeridos, etc.).
    \item \textbf{RF1.4}: Los cuadros de texto de las contraseñas deben tener la opción de mostrar y ocultar el contenido por motivos de seguridad.
  \end{itemize}
  
  \item \textbf{RF2}: Búsqueda y selección de revistas
  \begin{itemize}
    \item \textbf{RF2.1}: En la interfaz principal debe mostrar un botón que redirija a otra interfaz con una lista de revistas y su información relevante.
    \item \textbf{RF2.2}: La interfaz con la lista de revistas debe incluir una barra de búsqueda por texto y paginación para filtrar las revistas.
    \item \textbf{RF2.3}: La interfaz principal debe permitir a los usuarios seleccionar la categoría y el nombre de la revista utilizando listas \textit{dropdowns} con búsqueda por texto (activándose el botón para consultar el histórico del JCR de la revista).
    \item \textbf{RF2.4}: La interfaz principal debe permitir a los usuarios seleccionar uno o más modelos de predicción (activándose el botón de predicción del JCR).
    \item \textbf{RF2.5}: En la interfaz principal se deben validar los campos obligatorios y mostrar un mensaje de error en caso de que no se completen.
  \end{itemize}

  \newpage
  \item \textbf{RF3}: Interfaz de Historial JCR
  \begin{itemize}
    \item \textbf{RF3.1}: La interfaz de Historial JCR debe mostrar un gráfico con el JCR de la revista seleccionada en los últimos 5 años.
    \item \textbf{RF3.2}: Debe mostrar un gráfico con el cuartil de la revista seleccionada en los últimos 5 años.
  \end{itemize}
  
  \item \textbf{RF4}: Interfaz de Predicción JCR
  \begin{itemize}
    \item \textbf{RF4.1}: La interfaz de Predicción JCR debe permitir a los usuarios seleccionar uno o más modelos de predicción.
    \item \textbf{RF4.2}: Debe mostrar la predicción del JCR de cada modelo para el año actual.
    \item \textbf{RF4.3}: Debe mostrar un gráfico con la línea temporal del JCR de la revista y la predicción de los dos últimos años en diferentes colores para cada modelo.
    \item \textbf{RF4.4}: Los usuarios deben poder ocultar cualquier línea del gráfico haciendo clic sobre su leyenda.
  \end{itemize}
  
  \item \textbf{RF5}: Barra inferior y superior de las interfaces
  \begin{itemize}
    \item \textbf{RF5.1}: La barra inferior debe incluir enlaces a la política de privacidad y a los términos de uso de la aplicación.
    \item \textbf{RF5.2}: La barra inferior debe incluir el logo de la Universidad de Burgos que redirecciona a su página web oficial.
    \item \textbf{RF5.3}: La barra superior debe mostrar el logo de la aplicación que redirecciona siempre a la página principal.
    \item \textbf{RF5.4}: La barra superior debe incluir un menú \textit{dropdown} con opciones de internacionalización en español, inglés, italiano y francés.
    \item \textbf{RF5.5}: La barra superior debe mostrar un enlace a la ayuda en línea.
    \item \textbf{RF5.6}: La barra superior debe mostrar el nombre e imagen del usuario que ha iniciado sesión.
    \item \textbf{RF5.7}: El menú \textit{dropdown} del nombre del usuario debe permitir cerrar sesión (eliminando las variables de sesión y redireccionando a la interfaz de inicio de sesión) o editar el perfil (redireccionando a una página donde se pueda consultar y modificar la información del usuario).
  \end{itemize}
  
  \item \textbf{RF6}: Navegación
  \begin{itemize}
    \item \textbf{RF6.1}: Todas las interfaces deben tener navegación en forma de \textit{breadcrumbs} para indicar la ubicación actual del usuario en la aplicación.
  \end{itemize}
  
\end{itemize}

\subsection{Requisitos No Funcionales}

\begin{itemize}
  \item \textbf{RNF1}: Internacionalización
  \begin{itemize}
    \item \textbf{RNF1.1}: La aplicación debe tomar el idioma configurado en el navegador o el idioma <<más cercano>> disponible por defecto.
  \end{itemize}
  
  \item \textbf{RNF2}: Seguridad
  \begin{itemize}
    \item \textbf{RNF2.1}: Las contraseñas deben ser almacenadas de forma segura utilizando técnicas de cifrado adecuadas.
    \item \textbf{RNF2.2}: La aplicación debe protegerse contra ataques de fuerza bruta en los formularios de autenticación y registro.
  \end{itemize}
  
  \item \textbf{RNF3}: Rendimiento
  \begin{itemize}
    \item \textbf{RNF3.1}: La aplicación debe ser capaz de manejar de manera eficiente la carga de usuarios concurrentes.
    \item \textbf{RNF3.2}: Las consultas a la base de datos deben optimizarse para proporcionar respuestas rápidas a los usuarios.
  \end{itemize}
  
  \item \textbf{RNF4}: Usabilidad
  \begin{itemize}
    \item \textbf{RNF4.1}: La interfaz de usuario debe ser intuitiva y fácil de usar para los usuarios.
    \item \textbf{RNF4.2}: Los mensajes de error deben ser claros y descriptivos, brindando orientación sobre cómo solucionar los problemas.
  \end{itemize}
  
  \item \textbf{RNF5}: Mantenibilidad
  \begin{itemize}
    \item \textbf{RNF5.1}: El código fuente de la aplicación debe estar bien estructurado, modularizado y documentado.
    \item \textbf{RNF5.2}: El código debe seguir las mejores prácticas de desarrollo de \textit{software} y ser fácilmente mantenible y escalable.
  \end{itemize}
\end{itemize}


\section{Especificación de requisitos}
\label{sec:Especificación de requisitos}

En esta sección se enumerarán los distintos casos de uso para la aplicación web desarrollada.

\subsection{Actores}
Antes de comenzar con los casos de uso, se identifican los distintos actores que pueden interactuar con la aplicación y se definen a continuación:
\begin{itemize}
    \item \textbf{Usuario no registrado}: Referido a una persona que no ha creado una cuenta o perfil para entrar en la aplicación.
    \item \textbf{Usuario registrado}: Referido a una persona que ha creado una cuenta o perfil en la aplicación. Tendrá la posibilidad de modificar sus credenciales.
    \item \textbf{Administrador}: Usuario registrado con privilegios de administrador.
\end{itemize}

\subsection{Casos de Uso}
En cuanto a los casos de uso en si, se han identificado los siguientes:

\begin{itemize}
\item \textbf{CU-1} Autenticación (ver Tabla \ref{tab:cu1}).
\item \textbf{CU-2} Registro (ver Tabla \ref{tab:cu2}).
\item \textbf{CU-3} Interfaz Principal (ver Tabla \ref{tab:cu3}).
\item \textbf{CU-4} Interfaz de Lista de Revistas (ver Tabla \ref{tab:cu4}).
\item \textbf{CU-5} Buscar revistas (ver Tabla \ref{tab:cu5}).
\item \textbf{CU-6} Consultar historial de la revista (ver Tabla \ref{tab:cu6}).
\item \textbf{CU-7} Consultar predicción (ver Tabla \ref{tab:cu7}).
\item \textbf{CU-8} Consultar política de privacidad y términos de uso (ver Tabla \ref{tab:cu8}).
\item \textbf{CU-9} Redireccionar a la página de la Universidad de Burgos (ver Tabla \ref{tab:cu9}).
\item \textbf{CU-10} Redireccionar a la página principal (ver Tabla \ref{tab:cu10}).
\item \textbf{CU-11} Navegación con breadcrumbs (ver Tabla \ref{tab:cu11}).
\end{itemize}

\begin{table}[p]
\centering
\begin{tabularx}{\linewidth}{ p{0.21\columnwidth} p{0.71\columnwidth} }
\toprule
\textbf{CU-1} & \textbf{Autenticación}\\
\toprule
\textbf{Versión} & 1.0 \\
\textbf{Autor} & Gadea Lucas Pérez \\
\textbf{Requisitos asociados} & RF1, RF1.1, RF1.4  \\
\textbf{Descripción} & Este caso de uso describe el proceso de autenticación de los usuarios en la aplicación.\\
\textbf{Precondición} & El usuario debe estar registrado en la aplicación.\\
\textbf{Acciones} &
\begin{enumerate}
\def\labelenumi{\arabic{enumi}.}
\tightlist
\item El usuario ingresa su nombre de usuario y contraseña (pudiendo mostrar y ocultar su contenido).
\item El sistema verifica la validez de las credenciales ingresadas.
\item El sistema otorga acceso al usuario si las credenciales son válidas.
\end{enumerate}\\
\textbf{Postcondición} & El usuario ha iniciado sesión en la aplicación.\\
\textbf{Excepciones} &
\begin{itemize}
\item Si las credenciales son inválidas, se muestra un mensaje de error y se solicita al usuario que vuelva a ingresar las credenciales.
\item Si el usuario no está registrado, se muestra un mensaje indicando que debe registrarse antes de iniciar sesión.
\end{itemize}\\
\textbf{Importancia} & Alta\\
\bottomrule
\end{tabularx}
\caption{CU-1 Autenticación.}
\label{tab:cu1}
\end{table}


\begin{table}[p]
\centering
\begin{tabularx}{\linewidth}{ p{0.21\columnwidth} p{0.71\columnwidth} }
\toprule
\textbf{CU-2} & \textbf{Registro}\\
\toprule
\textbf{Versión} & 1.0 \\
\textbf{Autor} & Gadea Lucas Pérez \\
\textbf{Requisitos asociados} & RF1, RF1.3, RF1.4  \\
\textbf{Descripción} & Este caso de uso describe el proceso de registro de nuevos usuarios en la aplicación.\\
\textbf{Precondición} & - \\
\textbf{Acciones} &
\begin{enumerate}
\def\labelenumi{\arabic{enumi}.}
\tightlist
\item El usuario selecciona la opción de registro en la interfaz de inicio de sesión.
\item El usuario proporciona un nombre de usuario, una contraseña (pudiendo mostrar y ocultar su contenido), una dirección de correo electrónico válida y confirma la contraseña.
\item El sistema verifica que la contraseña tenga un formato válido y que la dirección de correo electrónico sea válida.
\item El sistema verifica que la contraseña y la confirmación de la contraseña coincidan.
\item El sistema crea un nuevo usuario en la base de datos con la información proporcionada.
\end{enumerate}\\
\textbf{Postcondición} & El usuario se ha registrado exitosamente.\\
\textbf{Excepciones} &
\begin{itemize}
\item Si la contraseña no cumple con el formato válido (i.e.: no tiene una longitud mínima de 8 caracteres, no contiene caracteres especiales requeridos, etc.), se muestra un mensaje de error y se solicita al usuario que ingrese una contraseña válida.
\item Si la dirección de correo electrónico no es válida, se muestra un mensaje de error y se solicita al usuario que ingrese una dirección de correo electrónico válida.
\item Si la contraseña y la confirmación de la contraseña no coinciden, se muestra un mensaje de error y se solicita al usuario que vuelva a ingresar la contraseña y su confirmación.
\end{itemize}\\
\textbf{Importancia} & Alta\\
\bottomrule
\end{tabularx}
\caption{CU-2 Registro.}
\label{tab:cu2}
\end{table}


\begin{table}[p]
\centering
\begin{tabularx}{\linewidth}{ p{0.21\columnwidth} p{0.71\columnwidth} }
\toprule
\textbf{CU-3} & \textbf{Interfaz Principal}\\
\toprule
\textbf{Versión} & 1.0 \\
\textbf{Autor} & Gadea Lucas Pérez \\
\textbf{Requisitos asociados} & RF2, RF2.1, RF2.2, RF2.3, RF2.4 \\
\textbf{Descripción} & Este caso de uso describe la interfaz principal de la aplicación, desde donde los usuarios pueden acceder a la lista de revistas y realizar búsquedas y filtrados.\\
\textbf{Precondición} & El usuario ha iniciado sesión en la aplicación. \\
\textbf{Acciones} &
\begin{enumerate}
\def\labelenumi{\arabic{enumi}.}
\tightlist
\item El usuario accede a la interfaz principal después de iniciar sesión correctamente.
\item La interfaz principal muestra un botón que redirige a la interfaz de lista de revistas.
\item La interfaz principal muestra listas desplegables (\textit{dropdowns}) con opciones de categorías y nombres de revistas.
\item El usuario puede utilizar la función de búsqueda por texto para filtrar las revistas en base a palabras clave.
\item La interfaz principal valida los campos obligatorios y muestra un mensaje de error si no se completan.
\end{enumerate}\\
\textbf{Postcondición} & - \\
\textbf{Excepciones} & - \\
\textbf{Importancia} & Alta\\
\bottomrule
\end{tabularx}
\caption{CU-3 Interfaz Principal.}
\label{tab:cu3}
\end{table}

\begin{table}[p]
\centering
\begin{tabularx}{\linewidth}{ p{0.21\columnwidth} p{0.71\columnwidth} }
\toprule
\textbf{CU-4} & \textbf{Interfaz de Lista de Revistas}\\
\toprule
\textbf{Versión} & 1.0 \\
\textbf{Autor} & Gadea Lucas Pérez \\
\textbf{Requisitos asociados} & RF2, RF2.2 \\
\textbf{Descripción} & Este caso de uso describe la interfaz que muestra una lista de revistas y su información relevante, con funcionalidades de búsqueda y paginación.\\
\textbf{Precondición} & El usuario ha accedido a la interfaz de lista de revistas desde la interfaz principal. \\
\textbf{Acciones} &
\begin{enumerate}
\def\labelenumi{\arabic{enumi}.}
\tightlist
\item La interfaz muestra una lista de revistas con su información relevante.
\item La interfaz incluye una barra de búsqueda por texto para filtrar las revistas en base a palabras clave.
\item La interfaz muestra una paginación para navegar entre las diferentes páginas de resultados.
\end{enumerate}\\
\textbf{Postcondición} & - \\
\textbf{Excepciones} & - \\
\textbf{Importancia} & Alta\\
\bottomrule
\end{tabularx}
\caption{CU-4 Interfaz de Lista de Revistas.}
\label{tab:cu4}
\end{table}


\begin{table}[p]
\centering
\begin{tabularx}{\linewidth}{ p{0.21\columnwidth} p{0.71\columnwidth} }
\toprule
\textbf{CU-5} & \textbf{Buscar revistas}\\
\toprule
\textbf{Versión} & 1.0 \\
\textbf{Autor} & Gadea Lucas Pérez \\
\textbf{Requisitos asociados} & RF2, RF2.2 \\
\textbf{Descripción} & Este caso de uso describe la funcionalidad para buscar revistas en la interfaz de lista de revistas desde la interfaz principal. Permite al usuario utilizar la barra de búsqueda por texto y navegar por la paginación de la lista de revistas. Además, el usuario puede acceder a la información adicional de una revista haciendo clic sobre su nombre. \\
\textbf{Precondición} & El usuario ha iniciado sesión en la aplicación y está en la interfaz principal. \\
\textbf{Acciones} &
\begin{enumerate}
\def\labelenumi{\arabic{enumi}.}
\tightlist
\item El usuario accede a la interfaz de lista de revistas desde la interfaz principal.
\item El usuario introduce un texto de búsqueda en la barra de búsqueda.
\item El usuario navega por la paginación de la lista de revistas.
\item El usuario hace clic sobre el nombre de una revista para acceder a su información adicional.
\end{enumerate}\\
\textbf{Postcondición} & - \\
\textbf{Excepciones} & - \\
\textbf{Importancia} & Alta\
\bottomrule
\end{tabularx}
\caption{CU-5 Buscar revistas.}
\label{tab:cu5}
\end{table}


\begin{table}[p]
\centering
\begin{tabularx}{\linewidth}{ p{0.21\columnwidth} p{0.71\columnwidth} }
\toprule
\textbf{CU-6} & \textbf{Consultar historial de la revista}\\
\toprule
\textbf{Versión} & 1.0 \\
\textbf{Autor} & Gadea Lucas Pérez \\
\textbf{Requisitos asociados} & RF2.3, RF3.1, RF3.2 \\
\textbf{Descripción} & Este caso de uso permite a los usuarios consultar el historial del JCR de una revista seleccionada. La interfaz principal debe permitir a los usuarios seleccionar la categoría y el nombre de la revista utilizando listas \textit{dropdowns} con búsqueda por texto, lo que activará el botón para consultar el histórico del JCR de la revista. \\
\textbf{Precondición} & El usuario ha iniciado sesión en la aplicación y está en la interfaz principal.\\
\textbf{Acciones} &
\begin{enumerate}
\def\labelenumi{\arabic{enumi}.}
\tightlist
\item El usuario selecciona la categoría de la revista en la lista \textit{dropdown}.
\item El usuario selecciona el nombre de la revista en la lista \textit{dropdown} o realiza una búsqueda por texto.
\item El usuario hace clic en el botón para consultar el histórico del JCR de la revista.
\item Se muestra la interfaz de Historial JCR con un gráfico del JCR de la revista en los últimos 5 años.
\item Se muestra un gráfico con el cuartil de la revista en los últimos 5 años.
\end{enumerate}\\
\textbf{Postcondición} & - \\
\textbf{Excepciones} & - \\
\textbf{Importancia} & Alta\\
\bottomrule
\end{tabularx}
\caption{CU-6 Consultar historial de la revista.}
\label{tab:cu6}
\end{table}


\begin{table}[p]
\centering
\begin{tabularx}{\linewidth}{ p{0.21\columnwidth} p{0.71\columnwidth} }
\toprule
\textbf{CU-7} & \textbf{Consultar predicción}\\
\toprule
\textbf{Versión} & 1.0 \\
\textbf{Autor} & Gadea Lucas Pérez \\
\textbf{Requisitos asociados} & RF2.4, RF4.1, RF4.2, RF4.3, RF4.4 \\
\textbf{Descripción} & Este caso de uso permite a los usuarios consultar la predicción del JCR de una revista utilizando uno o más modelos de predicción. La interfaz principal debe permitir a los usuarios seleccionar uno o más modelos de predicción, activando así el botón de predicción del JCR.\\
\textbf{Precondición} & El usuario ha iniciado sesión en la aplicación y está en la interfaz principal.\\
\textbf{Acciones} &
\begin{enumerate}
\def\labelenumi{\arabic{enumi}.}
\tightlist
\item El usuario selecciona uno o más modelos de predicción en la interfaz principal.
\item Se activa el botón de predicción del JCR.
\item El usuario hace clic en el botón de predicción del JCR.
\item Se muestra la interfaz de Predicción JCR.
\item El usuario selecciona uno o más modelos de predicción en la interfaz de Predicción JCR.
\item Se muestra la predicción del JCR de cada modelo para el año actual.
\item Se muestra un gráfico con la línea temporal del JCR de la revista y la predicción de los dos últimos años en diferentes colores para cada modelo.
\item El usuario puede ocultar cualquier línea del gráfico haciendo clic sobre su leyenda.
\end{enumerate}\\
\textbf{Postcondición} & - \\
\textbf{Excepciones} & - \\
\textbf{Importancia} & Alta\\
\bottomrule
\end{tabularx}
\caption{CU-7 Consultar predicción.}
\label{tab:cu7}
\end{table}


\begin{table}[p]
\centering
\begin{tabularx}{\linewidth}{ p{0.21\columnwidth} p{0.71\columnwidth} }
\toprule
\textbf{CU-8} & \textbf{Consultar política de privacidad y términos de uso}\\
\toprule
\textbf{Versión} & 1.0 \\
\textbf{Autor} & Gadea Lucas Pérez \\
\textbf{Requisitos asociados} & RF5.1 \\
\textbf{Descripción} & Este caso de uso muestra enlaces a la política de privacidad y a los términos de uso de la aplicación en la barra inferior de las interfaces.\\
\textbf{Precondición} & El usuario ha iniciado sesión en la aplicación.\\
\textbf{Acciones} &
\begin{enumerate}
\def\labelenumi{\arabic{enumi}.}
\tightlist
\item El usuario accede a cualquier interfaz donde se muestra la barra inferior.
\item El usuario selecciona el hiperenlace de política de privacidad o términos de uso.
\item El usuario es redireccionado a la interfaz correspondiente donde se muestra la información.
\end{enumerate}\\ 
\textbf{Postcondición} & - \\
\textbf{Excepciones} & - \\
\textbf{Importancia} & Media \\
\bottomrule
\end{tabularx}
\caption{CU-8: Consultar política de privacidad y términos de uso.}
\label{tab:cu8}
\end{table}

\begin{table}[p]
\centering
\begin{tabularx}{\linewidth}{ p{0.21\columnwidth} p{0.71\columnwidth} }
\toprule
\textbf{CU-9} & \textbf{Redireccionar a la página web oficial de la Universidad de Burgos}\\
\toprule
\textbf{Versión} & 1.0 \\
\textbf{Autor} & Gadea Lucas Pérez \\
\textbf{Requisitos asociados} & RF5.2 \\
\textbf{Descripción} & Este caso de uso redirige al usuario hacia la página web oficial de la Universidad de Burgos al hacer clic en su logo ubicado en la barra inferior.\\
\textbf{Precondición} & El usuario ha iniciado sesión en la aplicación.\\
\textbf{Acciones} &
\begin{enumerate}
\def\labelenumi{\arabic{enumi}.}
\tightlist
\item El usuario accede a cualquier interfaz donde se muestra la barra inferior.
\item El usuario selecciona el logo de la Universidad de Burgos.
\item El usuario es redireccionado a la página oficial de la universidad.
\end{enumerate}\\ 
\textbf{Postcondición} & - \\
\textbf{Excepciones} & - \\
\textbf{Importancia} & Baja \\
\bottomrule
\end{tabularx}
\caption{CU-9: Redireccionar a la página web oficial de la Universidad de Burgos.}
\label{tab:cu9}
\end{table}


\begin{table}[p]
\centering
\begin{tabularx}{\linewidth}{ p{0.21\columnwidth} p{0.71\columnwidth} }
\toprule
\textbf{CU-10} & \textbf{Redireccionar a la página principal}\\
\toprule
\textbf{Versión} & 1.0 \\
\textbf{Autor} & Gadea Lucas Pérez \\
\textbf{Requisitos asociados} & RF5.3 \\
\textbf{Descripción} & Este caso de uso redirige al usuario hacia la página principal de la aplicación al hacer clic en el logo de la aplicación ubicado en la barra superior.\\
\textbf{Precondición} & El usuario ha iniciado sesión en la aplicación.\\
\textbf{Acciones} &
\begin{enumerate}
\def\labelenumi{\arabic{enumi}.}
\tightlist
\item El usuario accede a cualquier interfaz donde se muestra la barra superior.
\item El usuario selecciona el logo de la aplicación.
\item El usuario es redireccionado a la página principal de la aplicación.
\end{enumerate}\\ 
\textbf{Postcondición} & - \\
\textbf{Excepciones} & - \\
\textbf{Importancia} & Alta \\
\bottomrule
\end{tabularx}
\caption{CU-10: Redireccionar a la página principal.}
\label{tab:cu10}
\end{table}


\begin{table}[p]
\centering
\begin{tabularx}{\linewidth}{ p{0.21\columnwidth} p{0.71\columnwidth} }
\toprule
\textbf{CU-11} & \textbf{Navegación con breadcrumbs}\\
\toprule
\textbf{Versión} & 1.0 \\
\textbf{Autor} & Gadea Lucas Pérez \\
\textbf{Requisitos asociados} & RF6.1 \\
\textbf{Descripción} & Este caso de uso abarca la implementación de la navegación en forma de breadcrumbs en todas las interfaces de la aplicación. Los breadcrumbs permiten al usuario conocer la ubicación actual dentro de la aplicación y navegar hacia atrás de manera intuitiva.\\
\textbf{Precondición} & El usuario ha iniciado sesión en la aplicación.\\
\textbf{Acciones} &
\begin{enumerate}
\def\labelenumi{\arabic{enumi}.}
\tightlist
\item El usuario accede a cualquier interfaz donde se muestran las breadcrums.
\item El usuario pulsa sobre el hipervínculo del \textit{breadcrumb}.
\item El usuario es redireccionado a la página correspondiente.
\end{enumerate}\\ 
\textbf{Postcondición} & - \\
\textbf{Excepciones} & - \\
\textbf{Importancia} & Media \\
\bottomrule
\end{tabularx}
\caption{CU-11: Navegación y breadcrumbs.}
\label{tab:cu11}
\end{table}



\subsection{Diagrama de casos de usos}

Finalmente, para poder visualizar todos estos casos, se incluye un diagrama de casos de uso (ver Figura \ref{fig:cu}).

\imagen{cu}{Diagrama de Casos de Uso}{1.1}
\apendice{Especificación de diseño}

\section{Introducción}

El objetivo principal de esta especificación de diseño es brindar una guía clara y completa para el desarrollo de la aplicación web, proporcionando una base sólida sobre la cual los desarrolladores puedan implementar el sistema de manera efectiva. Asimismo, este informe servirá como referencia y documentación para futuras etapas del proyecto y posibles mejoras o expansiones.

Así pues, en esta sección se describirán y analizarán en detalle los diferentes aspectos del diseño de la aplicación web, abarcando desde la estructura y organización de los datos hasta los procedimientos y algoritmos que permiten su funcionamiento adecuado. Además, se profundizará en el diseño arquitectónico, que define la distribución de los componentes y la interacción entre ellos.

\section{Diseño de los datos}

Comenzaremos detallando el diseño de los datos de la aplicación. En este caso, para el almacenamiento de la información, se ha utilizado una base de datos PostgreSQL, la cual es accedida a través de la aplicación haciendo uso de la librería de Python psycopg2. 

\subsection{Diccionario de datos}
A continuación, se describirán las tablas que componen la base de datos y se explicará la estructura y los atributos de cada una de ellas.

\subsubsection{Tabla <<Users>>}

Esta tabla almacena la información de los usuarios registrados en la aplicación. Los atributos de la tabla son los siguientes:
\begin{itemize}
    \item \texttt{username}: campo de tipo \texttt{VARCHAR(255)} que guarda el nombre de usuario.
    \item \texttt{password}: campo de tipo \texttt{VARCHAR(255)} que almacena la contraseña del usuario.
    \item \texttt{email}: campo de tipo \texttt{VARCHAR(255)} que sirve como clave primaria y guarda la dirección de correo electrónico del usuario.
    \item \texttt{admin}: campo de tipo \texttt{BOOLEAN} que indica si el usuario tiene privilegios de administrador.
\end{itemize}

\subsubsection{Tabla <<Modelos>>}

En esta tabla se registran los modelos utilizados en la aplicación. Los atributos de la tabla son los siguientes:
\begin{itemize}
    \item \texttt{id}: campo de tipo \texttt{SERIAL}  que actúa como clave primaria y asigna un identificador único a cada modelo.
    \item \texttt{nombre}: campo de tipo \texttt{TEXT} que almacena el nombre del modelo.
    \item \texttt{rmse}: campo de tipo \texttt{FLOAT} que guarda el valor del error cuadrático medio asociado al modelo.
\end{itemize}


\subsubsection{Tabla <<Revista>>}

Esta tabla contiene la información de las revistas presentes en la aplicación. Los atributos de la tabla son los siguientes:
\begin{itemize}
    \item \texttt{nombre}: campo de tipo \texttt{VARCHAR(255)} que sirve como clave primaria y almacena el nombre de la revista.
    \item \texttt{ISSN}: campo de tipo \texttt{VARCHAR(9)} que garantiza la unicidad y no nulidad del International Standard Serial Number de la revista.
    \item \texttt{categoria}: campo de tipo \texttt{VARCHAR(255)} que guarda la categoría a la que pertenece la revista.
\end{itemize}

\subsubsection{Tabla <<Revista\_jcr>>}

Esta tabla registra las métricas JCR (Journal Citation Reports) asociadas a las revistas. Los atributos de la tabla son los siguientes:

\begin{itemize}
    \item \texttt{nombre}: campo de tipo \texttt{VARCHAR(255)} que guarda el nombre de la revista.
    \item \texttt{fecha}: campo de tipo \texttt{NUMERIC} que indica la fecha de la métrica JCR.
    \item \texttt{jcr}: campo de tipo \texttt{FLOAT} que almacena el valor de la métrica JCR para la revista en la fecha correspondiente.
    \item \texttt{citas}: campo de tipo \texttt{NUMERIC} que indica la cantidad de citas recibidas por la revista en la fecha dada.
    \item \texttt{diff}: campo de tipo \texttt{FLOAT} que representa la diferencia en la métrica JCR real en comparación con el JCR calculado por el alumno (en la fecha dada).
    \item \texttt{cuartil}: campo de tipo \texttt{VARCHAR(3)} que indica el cuartil al que ha pertenecido la revista en la fecha dada.
\end{itemize}
 
Estas tablas y sus respectivos atributos han sido diseñados de manera coherente para almacenar y organizar la información necesaria en la aplicación web de una forma sencilla.

\subsection{Optimización de la base de datos}

Para mejorar el rendimiento y la eficiencia de las consultas en la base de datos, se han aplicado diferentes técnicas de optimización. Una de estas técnicas es la creación de índices, que permiten acceder rápidamente a los registros de una tabla en función de los valores de una columna o conjunto de columnas. En el caso de la tabla <<Revista>>, se ha creado un índice llamado \texttt{nombre\_index} en la columna \texttt{nombre} mediante la instrucción \texttt{CREATE INDEX nombre\_index ON revista (nombre)}. Esto facilita la búsqueda y recuperación de registros en la tabla <<Revista>> basándose en el valor de la columna \texttt{nombre}.

Además de crear índices, se ha utilizado el comando \texttt{ANALYZE} para actualizar las estadísticas de las tablas. Esto permite al sistema de gestión de bases de datos (SGBD) recopilar información actualizada sobre la distribución de los datos y la estructura de las tablas. En el caso de la tabla <<Revista>>, se ha aplicado el comando \texttt{ANALYZE revista} para asegurar que las estadísticas se encuentren actualizadas y que el SGBD pueda tomar decisiones más acertadas al ejecutar consultas sobre esta tabla. Del mismo modo, se ha realizado el análisis del resto de tablas de la base de datos.

Estas técnicas permiten acelerar las consultas y mejorar la eficiencia en la recuperación de datos, lo que resulta en una aplicación web más rápida y con mejor respuesta para los usuarios.

\subsection{Diagramas de datos}
A continuación se mostrará un sencillo diagrama relacional (ver Figura \ref{fig:rel}).

% Diagrama relacional
\imagen{rel}{Diagrama Relacional}{1}


\section{Diseño procedimental}
TODO: diagramas de secuencia

\section{Diseño arquitectónico}

Se pretende desarrollar una aplicación web sencilla que utilizará una arquitectura cliente-servidor de dos capas.

La funcionalidad principal de la aplicación será predecir el índice de impacto de una lista de revistas elegida por el usuario. Además, se proporcionarán gráficos e información adicional para ayudar a los usuarios a entender mejor los resultados. Por otro lado, se podrán distinguir dos perfiles de usuario: el usuario normal y el administrador. El administrador tendrá permisos para ``reentrenar'' la red que predice los índices y gestionar al resto de usuarios.

El \textit{backend} se programará en Python y el \textit{frontend} con HTML, CSS y JavaScript. Para ello, se utilizará Flask, que se trata de un \textit{framework} de Python para el desarrollo de aplicaciones web.

Para la conexión con la base de datos, se hará uso de \textit{psycopg2}, que es un módulo de Python que proporciona una interfaz para conectarse y interactuar con bases de datos PostgreSQL. 
Es una de las librerías más populares y ampliamente utilizadas para trabajar con PostgreSQL en Python. Además, es compatible con la mayoría de las versiones de Python y es muy fácil de utilizar.

Por otro lado, para la autenticación, se han evaluado varias opciones que ofrece Flask, a saber: Flask-Login, Flask-Security y Flask-JWT. 
Finalmente, se ha elegido Flask-Login debido a su sencillez y facilidad de uso. Esta opción permite al usuario iniciar sesión con un nombre de usuario y una contraseña y almacena la información de la sesión en las \textit{cookies} del navegador. Este paquete incluye la protección de rutas con decoradores.


\section{Diseño de Interfaces}
La aplicación estará conformada por seis interfaces principales, que se detallan a continuación:

\begin{itemize}
    \item Interfaces iniciales con funciones de inicio de sesión y registro, verificando la validez de los datos y permitiendo recuperar contraseñas.
    \item Interfaz principal para seleccionar una revista y un año, con opción de calcular el JCR y elegir modelos de predicción.
    \item Interfaces con el histórico y la predicción del JCR, además de gráficos representativos.
    \item Interfaz con una lista de las revistas disponibles con barra de búsqueda.
    \item Área de usuario donde modificar las credenciales.
\end{itemize}

% METER PROTOTIPOS 





\apendice{Documentación técnica de programación}

\section{Introducción}



\section{Estructura de directorios}

\section{Manual del programador}

Esta sección está destinada a proporcionar información detallada sobre cómo utilizar el programa desarrollado en el proyecto. Aquí se describen las funciones y características clave de dicho software, así como los detalles sobre la configuración y la integración del programa con otros sistemas.

\subsection{Base de datos}
Para aquellos usuarios de sistemas operativos Windows, es importante tener en cuenta un detalle en el momento de obtener información de la base de datos en PostgresSQL. Cuando se recolectan los datos en formato CSV, es necesario cambiar la ruta de generación de estos archivos a la carpeta <<Public>> \footnote{Directorio que permite compartir archivos entre distintos usuarios en un mismo sistema Windows. Microsoft ha mantenido este directorio desde la versión de WindowsXP.}. Esto se debe a que el usuario por defecto que genera PostgresSQL (usuario <<postgres>>), no tiene los permisos suficientes para leer e importar datos de los CSV en cuestión ya que, para ello, se utiliza el comando \textit{copy} (que solo puede ser ejecutado por <<Superusers>>)\cite{Dominguez2020}. Sin embargo, la carpeta <<Public>> cuenta con los permisos necesarios para que el usuario que crea Postgres pueda acceder y leer los datos. Por lo tanto, es importante seguir este paso para garantizar que se pueda obtener correctamente la información recogida previamente.


\section{Compilación, instalación y ejecución del proyecto}

\section{Pruebas del sistema}

\apendice{Documentación de usuario}

\section{Introducción}

\section{Requisitos de usuarios}

\section{Instalación}

\section{Manual del usuario}


\subsection{Entorno virtualizado}

Se trata de un entorno aislado que te permite instalar y gestionar de forma independiente las bibliotecas y paquetes de Python que se necesitan para el proyecto, sin afectar a otros proyectos o al sistema operativo. Esto nos permite evitar  problemas de dependencias y compatibilidad entre diferentes versiones de bibliotecas, además de facilitar la gestión y el mantenimiento del proyectos.
 
Miniconda, por su parte, es una distribución de Python que incluye un gestor de paquetes llamado Conda, que te permite instalar y gestionar de forma sencilla bibliotecas y paquetes de Python.

Así pues, se procede a crear un entorno virtualizado de Python utilizando Miniconda. El entorno en cuestión se encuentra empaquetado
en la carpeta comprimida "environment.zip". Para poder hacer uso del
entorno bastará con ejecutar el siguiente comando:
\texttt{conda activate /ruta/al/environment.zip}



\apendice{Proyectos de pruebas}

\section{Introducción}
En el desarrollo de \textit{software}, es esencial contar con una metodología adecuada para garantizar la calidad del producto final. En este sentido, el Test Driven Development (TDD) se ha convertido en una de las metodologías más populares en la actualidad. TDD consiste en escribir las pruebas antes de comenzar a desarrollar el código, lo que permite detectar errores de manera temprana y asegurar que el \textit{software} cumpla con los requisitos establecidos. En este apartado, presentaremos los proyectos de pruebas que hemos llevado a cabo utilizando TDD como metodología principal.

Esta metodología se divide en tres fases principales: la fase roja, la fase verde y la fase morada. 
\begin{itemize}
    \item En la fase roja, se escriben las pruebas unitarias para el código que se desea implementar, las cuales deben fallar inicialmente (puesto que aún no existe el código a probar). 
    \item En la fase verde, se escribe el código mínimo necesario para que las pruebas pasen satisfactoriamente.
    \item Finalmente, en la fase morada, se realiza la refactorización del código escrito anteriormente para mejorar su calidad y mantenerlo limpio y legible.
\end{itemize}

 Este enfoque iterativo permite a los desarrolladores garantizar la calidad y funcionalidad del software en cada paso del proceso de desarrollo y ayuda a reducir los errores y \textit{bugs} que puedan surgir en el proceso. 

\imagen{fasesTDD}{Fases del la metodología TDD}{0.4}

\section{Prototipo de Crossref}
Este es el prototipo final para la extracción de los datos bibliográficos. A continuación, se documentará la información más relevante del proyecto de pruebas llevado a cabo.

\subsection{Diseño y arquitectura}
Puesto que el \textit{software} a probar está compuesto de un solo \textit{script} de Python, se ha generado un único módulo de pruebas llamado \texttt{crossref\_test.py}. 

\subsection{Fase Roja}
A lo largo de esta fase se definen las pruebas necesarias para crear un código que extraiga correctamente los datos de Crossref. Los casos de prueba diseñados son los siguientes:

NOTA -> Ver pregunta sobre la tabla al final del documento
	
\tablatres{Casos de prueba para el prototipo de Crossref}{|l|l|l|l|}{4}{casosdepruebacrossref}{\textbf{ID} & \textbf{Nombre} & \textbf{Descipción} & \parbox{55}{\textbf{Salida esperada}}\\}{
C01 & \parbox{60}{Revista existente.} & 
\parbox{200}{\vspace{0.4cm}Método test\_journal\_exist. Se trata de extraer una revista que previamente se conoce que existe.\vspace{0.4cm}} & \parbox{55}{True.}\\
C02 & \parbox{60}{Revista inexistente.} & 
\parbox{200}{\vspace{0.4cm} Método test\_journal\_exist. Se trata de extraer una revista que previamente se conoce que no existe.\vspace{0.4cm}} &\parbox{55}{False.}\\
C03 & \parbox{60}{Cargar datos de entrada.} & 
\parbox{200}{\vspace{0.4cm}Método test\_load\_data. Se comprueba que se cargan bien los datos del CSV de entrada.\vspace{0.4cm}}&\parbox{55}{Carga correcta.}\\
C04 & \parbox{60}{Formateado I.} & 
\parbox{200}{\vspace{0.4cm}Método test\_formatData. Se comprueba que se eliminan bien las etiquetas HTML y los espacios.\vspace{0.4cm}}&\parbox{55}{Cadena correcta.}\\
C05 & \parbox{60}{Formateado II.} & 
\parbox{200}{\vspace{0.4cm}Método test\_formatData. Se comprueba que se eliminan bien las etiquetas "href".\vspace{0.4cm}}& \parbox{55}{Cadena correcta.}\\
C06 & \parbox{60}{Formateado III.} &
\parbox{200}{\vspace{0.4cm}Método test\_formatData. Se comprueba que una cadena correcta no cambia.\vspace{0.4cm}} & \parbox{55}{Cadena sin modificaciones.}\\
C07 & \parbox{60}{Formateado IV.} & 
\parbox{200}{\vspace{0.4cm}Método test\_formatData. Se comprueba que no hay campos vacíos (rellenar con "n.m.").\vspace{0.4cm}} & \parbox{55}{Cadena sin valores vacíos.}\\
C08 & \parbox{60}{Formato de salida I.} & 
\parbox{200}{\vspace{0.4cm}Método test\_getArticles\_returns\_list\_
of\_dicts. Se comprueba que se obtiene una lista de artículos.\vspace{0.4cm}} & \parbox{55}{Lista de diccionarios.}\\
C09 & \parbox{60}{Formato de salida II.} & 
\parbox{200}{\vspace{0.4cm}Método test\_getArticles\_returns\_valid
\_keys. Se comprueba que las claves del diccionario son correctas.\vspace{0.4cm}} & \parbox{55}{Claves correctas.}\\
C10 & \parbox{60}{Número de resultados.} & 
\parbox{200}{\vspace{0.4cm}Método test\_getArticles\_correct\_
article\_count. Se comprueba que se obtiene el número correcto de artículos que aparecen en Crossref.\vspace{0.4cm}} & \parbox{55}{Número correcto.}\\
}

\tabla{Casos de prueba para el prototipo de Crossref}{|l|l|l|l|}{4}{casosdepruebacrossref}{\textbf{ID} & \textbf{Nombre} & \textbf{Descipción} & \parbox{55}{\textbf{Salida esperada}}\\}{
C11 & \parbox{60}{Excepción de ISSN.} & 
\parbox{200}{\vspace{0.4cm}Método test\_getArticles\_invalid\_issn\_
raises\_error. Se intenta buscar un ISSN incorrecto.\vspace{0.4cm}} & \parbox{55}{False.}\\
C12 & \parbox{60}{Excepción de año I.} & 
\parbox{200}{\vspace{0.4cm}Método test\_getArticles\_invalid\_year. Se introducen los años de búsqueda en el orden incorrecto.\vspace{0.4cm}} &\parbox{55}{False.}\\
C13 & \parbox{60}{Excepción de año II.} & 
\parbox{200}{\vspace{0.4cm}Método test\_getArticles\_invalid\_year. Se introducen años inválidos.\vspace{0.4cm}} &\parbox{55}{False.}\\
C14 & \parbox{60}{Ejecución completa I.} & 
\parbox{200}{\vspace{0.4cm}Método test\_main. Pre-requisito: CSV de entrada de datos. Se comprueba que se genera una gráfica al finalizar el proceso. Post-requisito: Imagen con la gráfica generado.\vspace{0.4cm}} &\parbox{55}{True.}\\
C15 & \parbox{60}{Ejecución completa II.} & 
\parbox{200}{\vspace{0.4cm}Método test\_main. Precondición: CSV de entrada de datos. Se comprueba que se generan los ficheros CSVs de resultados al finalizar el proceso. Post-requisito: Ficheros CSV generados.\vspace{0.4cm}}&\parbox{55}{True.}\\
C16 & \parbox{60}{Ejecución completa III.} & 
\parbox{200}{\vspace{0.4cm}Método test\_main. Precondición: CSV de entrada de datos. Se comprueba que se genera el fichero de log al finalizar el proceso. Postcondición: Ficheros de log generados.\vspace{0.4cm}}&\parbox{55}{True.}\\
C17 & \parbox{60}{Ejecución completa IV.} & 
\parbox{200}{\vspace{0.4cm}Método test\_main. Precondición: CSV de entrada de datos y fichero de log generado correctamente. Se comprueba que la información de log es correcta.\vspace{0.4cm}}&\parbox{55}{True.}\\
}

\subsection{Fase Verde}
Tras la creación del código, finalmente se consigue una versión que pasa todos los test programados en la fase anterior, como se puede observar en la siguiente ilustración:

\imagen{pruebasPasadas}{Pruebas pasadas exitosamente}{0.7}

\subsection{Fase Morada}
En esta fase se han realizado acciones de refactorización. A saber:

\begin{itemize}
    \item Documentación de código:
    Se han incluido las cabeceras y comentarios más importantes en el código para facilitar su comprensión y claridad. Asimismo, se han eliminado comentarios redundantes o innecesarios.

    \item Ordenación de \textit{imports}:
    Se han ordenado los \textit{imports} del código de forma que se identifique fácilmente la naturaleza de los paquetes y módulos empleados. De esta forma, se agrupan al inicio los módulos integrados de Python, posteriormente los paquetes de terceros y, finalmente, módulos o paquetes locales.
    
    \item Mensajes de salida I:
    En lugar de imprimir mensajes de salida en la consola, (ant. \verb|print("[ERROR] Se ha producido un error")|), ahora se utilizan registros de \textit{logging} para registrar los eventos de la función (por ejemplo, \verb|logging.error('Se ha producido un error')|). Los registros son una forma más eficiente y flexible de registrar información que imprimir en la consola.
        
    \item Mensajes de salida II:
    Utiliza f-strings en lugar de concatenar cadenas para mejorar el rendimiento y la legibilidad del código. Ejemplo: \verb|logging.error('Se ha producido un error'+ codigo_error)| por \verb|logging.error(f'Se ha producido un error {codigo_error}')|.

    \item Acceso a diccionarios I:
    En lugar de especificar los parámetros de búsqueda utilizando una cadena de consulta, se puede usar un diccionario y pasarlos como argumentos de palabras clave al método \verb|requests.get()|. Esto puede hacer que el código sea más legible.

    \item Acceso a diccionarios II:
    En lugar de crear un diccionario vacío "author" y luego verificar si existe una lista de autores, se puede usar el método \verb|get()| directamente en la lista de autores para obtener el primer elemento de la lista.

    \item Lista de comprensión:
    En lugar de crear un diccionario temporal publicacion para cada artículo, se puede usar una lista de comprensión para crear una lista de diccionarios de artículos. Esto puede hacer que el código sea más conciso.
	
    
\end{itemize}





\section{PREGUNTA SOBRE TABLA}

He hecho una tabla general con los casos de prueba muy resumidos... ¿crees que es suficiente, o sería conveniente elaborar una tabla para cada uno de los casos de prueba (como la de abajo)?

\begin{table}[p]
	\centering
	\begin{tabularx}{\linewidth}{ p{0.21\columnwidth} p{0.71\columnwidth} }
		\toprule
		\textbf{CP-1}    & \textbf{Ejemplo de caso de prueba}\\
		\toprule
		\textbf{Versión}              & 1.0    \\
		\textbf{Autor}                & Alumno \\
		\textbf{Requisitos asociados} & RF-xx, RF-xx \\
		\textbf{Descripción}          & La descripción del CP \\
		\textbf{Precondición}         & Precondiciones (podría haber más de una) \\
		\textbf{Acciones}             &
		\begin{enumerate}
			\def\labelenumi{\arabic{enumi}.}
			\tightlist
			\item Pasos del CP
			\item Pasos del CP (añadir tantos como sean necesarios)
		\end{enumerate}\\
		\textbf{Postcondición}        & Postcondiciones (podría haber más de una) \\
		\textbf{Excepciones}          & Excepciones \\
		\textbf{Importancia}          & Alta o Media o Baja... \\
		\bottomrule
	\end{tabularx}
	\caption{CP-1 Nombre del caso de prueba.}
\end{table}


	
	


\bibliographystyle{plain}
\bibliography{bibliografiaAnexos}

\end{document}
